% !TEX encoding = UTF-8
\documentclass[UTF8]{ctexart}

\usepackage[utf8]{inputenc}
\usepackage{graphicx}
\usepackage{geometry}
\geometry{a4paper}
\geometry{left=2.5cm,right=2.5cm,top=2.8cm,bottom=1.3cm}

\usepackage{booktabs}
\usepackage{array}
\usepackage{paralist}
\usepackage{verbatim}
\usepackage{subfig}
\usepackage{amsmath}
\usepackage{mathtools}
\usepackage{listings}
\usepackage[table]{xcolor}
\usepackage{lastpage}
\usepackage{url}

% Using hyperref for improved ref character
\usepackage[colorlinks,linkcolor=black,anchorcolor=black,
citecolor=black,CJKbookmarks=True]{hyperref}

% For picture drawing
\usepackage[all]{xy}

% For code inserting. Set features.
\lstset{
alsolanguage=matlab,
tabsize=4,
keepspaces=true,
numbers=left,
numberstyle=\tiny,
keywordstyle=\color{blue!70} \bfseries,
commentstyle=\color{red!50!green!50!blue!50},
frame=shadowbox,
breaklines,
showspaces=false,
showstringspaces=false,
showtabs=false,
rulesepcolor=\color{red!20!green!20!blue!20},
extendedchars=false,
escapeinside=``
}

% Set the font of page header
\usepackage{fancyhdr}
\pagestyle{fancy}
\lhead{传感器特性实验报告}
\chead{}
\rhead{Page \thepage/\pageref{LastPage}}
\cfoot{}
\rfoot{}
\lfoot{}

\usepackage{sectsty}

\usepackage[nottoc]{tocbibind}
\usepackage[titles,subfigure]{tocloft}
\renewcommand{\cftsecfont}{\rmfamily\mdseries\upshape}
\renewcommand{\cftsecpagefont}{\rmfamily\mdseries\upshape}

% Set number of ref to be relevent to section number
\renewcommand{\theequation}{\arabic{section}.\arabic{equation}}
\renewcommand{\thefigure}{\arabic{section}-\arabic{figure}}
\renewcommand{\thetable}{\arabic{section}-\arabic{table}}
\makeatletter
\@addtoreset{equation}{section}
\@addtoreset{figure}{section}
\@addtoreset{table}{section}
\makeatother

% Set the font of the reference
\bibliographystyle{unsrt}

% Define user\rq{}s color
\usepackage{colortbl}
\definecolor{lightgray}{gray}{.9}
\definecolor{thickgray}{gray}{.6}

\usepackage{multirow}

% 首行缩进
\usepackage{indentfirst}

% Set section numbering
\CTEXsetup[number={}]{part}
\renewcommand{\thepart}{}
\usepackage{titlesec}
\titleformat{\part}[block]{\color{blue}\huge\bfseries\filcenter}{}{1em}{}

%\usepackage{ulem}
%\usepackage{indentfirst}
%\setlength\textwidth{300.0pt}
%

% 重定义字体命令
\newcommand{\song}{\CJKfamily{song}}    % 宋体   (Windows自带simsun.ttf)
\newcommand{\fs}{\CJKfamily{fs}}        % 仿宋体 (华天字库htfs.ttf)
\newcommand{\kai}{\CJKfamily{kai}}      % 楷体   (华天字库htkai.ttf)
\newcommand{\hei}{\CJKfamily{hei}}      % 黑体   (Windows自带simhei.ttf)
\newcommand{\li}{\CJKfamily{li}}        % 隶书   (Windows自带simli.ttf)
\newcommand{\you}{\CJKfamily{you}}      % 幼圆体 (Windows自带simyou.ttf)
%%%  以上六种字体均为标准 GBK 字体, 包括 GBK 繁体字和一些不常用字, 推荐!!!

\newcommand{\xs}{\CJKfamily{xs}}
\newcommand{\shu}{\CJKfamily{shu}}      % 舒体   (方正字库fzstk.ttf)
%  \newcommand{\yourcommand}[参数个数]{内容}   [参数个数]这个是可选的。
%  例如  \newcommand{\you}{\CJKfamily{you}}  用\you 来代替 \CJKfamily{you} ，少输入很多字哦。
%字号设置
\newcommand{\chuhao}{\fontsize{42pt}{\baselineskip}\selectfont}
\newcommand{\xiaochuhao}{\fontsize{36pt}{\baselineskip}\selectfont}
\newcommand{\yihao}{\fontsize{28pt}{\baselineskip}\selectfont}
\newcommand{\xiaoyihao}{\fontsize{24pt}{\baselineskip}\selectfont}
\newcommand{\erhao}{\fontsize{21pt}{\baselineskip}\selectfont}
\newcommand{\xiaoerhao}{\fontsize{18pt}{\baselineskip}\selectfont}
\newcommand{\sanhao}{\fontsize{15.75pt}{\baselineskip}\selectfont}
\newcommand{\sihao}{\fontsize{14pt}{\baselineskip}\selectfont}
\newcommand{\xiaosihao}{\fontsize{12pt}{\baselineskip}\selectfont}
\newcommand{\wuhao}{\fontsize{10.5pt}{\baselineskip}\selectfont}
\newcommand{\xiaowuhao}{\fontsize{9pt}{\baselineskip}\selectfont}
\newcommand{\liuhao}{\fontsize{7.875pt}{\baselineskip}\selectfont}
\newcommand{\qihao}{\fontsize{5.25pt}{\baselineskip}\selectfont}
% \baselineskip | distance from bottom of one line of a paragraph to bottom of the next line.  基本行距
%  只有使用\selectfont命令之后，\fontzize{}{}的设置才能生效
%  可以用数字表示{11pt}:单倍行距

\begin{document}
%%%%%%%%%%%%%%%%%%%%%%%%%%%%封面与目录%%%%%%%%%%%%%%%%%%%%%%%%%%%%%%
\begin{titlepage}
\begin{center}
% Upper part of the page
\includegraphics[width=0.25\textwidth]{resource/logo.jpg}\\[1cm]
\textsc{\LARGE Department of Automation}\\[1.5cm]
\fs{\Large 模式识别基础第二次作业}\\[0.5cm]
% Title
\hrulefill
\\[0.8cm]{\centering \huge \hei 用身高体重数据进行性别分类的实验（二）}\\[0.4cm]
\hrulefill
\\[4cm]

% Author and supervisor
\begin{tabbing}       %tabbing  列表

 \hspace*{5cm} \= \hspace{2.6cm} \= \kill
 % \=     in tabbing environment, sets a tab stop
 % \kill  in a\tabbing environment, deletes previous line so tabs can be set without outputting text.
 % \>     in tabbing environment is a forward tab.

\>{\fs\sihao\textbf {班\hspace{1cm}级 \ \ ：}}\>  {\centering\fs\sihao\textbf{~~~~~~~~~自~~3~2}} \\
\\
\>{\fs\sihao\textbf {姓\hspace{1cm}名 \ \ ：}}\>  {\centering\fs\sihao\textbf{~~~~~~~~陈~昊~楠}}\\
\\
\>{\fs\sihao\textbf {学\hspace{1cm}号 \ \ ：}}\>  {\centering\fs\sihao\textbf{~~~~~~2013011449}}\\
\\
\>{\fs\sihao\textbf {授课教师 \ \ ：}}\>  {\centering\fs\sihao\textbf{~~~~~~~~张~学~工}} \\

\end{tabbing}
\vfill
{\large \today}
\end{center}
\end{titlepage}

\tableofcontents
\clearpage

%%%%%%%%%%%%%%%%%%%%%%%%%%正文部分%%%%%%%%%%%%%%%%%%%%%%%%%%%%%%%%%%

\section{实验内容及结果}
\begin{enumerate}
	\item 用 {\ttfamily dataset3} 作为训练数据，用 {\ttfamily dataset4} 作为测试数据，采用不同的特征、训练样本数、分类方法进行比较实验，观察、分析实验结果的异同。
	\item 
\end{enumerate}

\subsection{Bayes最小错误率分类}
\paragraph{}
\paragraph{实验结果} 
\subsection{}

\section{分析与体会}


\section{数据与方法}
\subsection{随机采样方法}
采用{\ttfamily sample}函数生成一定范围内不重复的随机整数。例如：
\begin{lstlisting}
	sample(1:size(df_train, 1), 10, replace = false)
\end{lstlisting}
函数具体说明参见 \url{http://statsbasejl.readthedocs.io/en/latest/sampling.html}。挑选出的样本参见源程序 \url{http://nbviewer.jupyter.org/github/chaonan99/PR_hw/blob/master/h2/src/h2.ipynb}。

\subsection{程序说明}
程序采用{\ttfamily Julia}语言，由{\ttfamily IJulia}提供对{\ttfamily Jupyter Notebook}的支持。
\paragraph{简要方法描述} 见实验内容部分。
\paragraph{程序来源} 源代码均为本人编写。用到的包参见源程序中的说明。
\paragraph{参数说明} 参见源程序中的说明。

\begin{figure}[htbp]
\centering
\includegraphics[width=11cm]{resource/logo.jpg}
\caption{全桥测试电路}
\label{fig:fullarm}
\end{figure}

\end{document}
